\documentclass[11pt,a4paper]{ltjsarticle}
\usepackage{amsmath}
\begin{document}

\section{和文の段落}

日本語の本文がそのまま組まれる段落である。行分割は禁則処理を含み、
約物の詰めと欧文との混植が同時に起きるので、シェーピングと
JFMグルーの両方がこの一段落だけで検査できる。English words mixed in
run alongside 和文 with inter-script glue.

\section{数式との混植}

和文の途中に $\alpha + \beta = \gamma$ のような数式が入り、別行立ての
式も続く:
\begin{equation}
  f(x) = \frac{1}{\sqrt{2\pi}} e^{-x^2/2}
\end{equation}

式のあとの段落。ここは前の式との間隔（\verb|\abovedisplayskip| 系）が
連続ランと一致するかを見るための本文である。長めに書いて二行以上に
折り返させる。

\section{箇条書き}

\begin{enumerate}
  \item 一つ目の項目。短い。
  \item 二つ目の項目。こちらは折り返しが起きる程度に長く書いてあり、
        字下げの継続も確認できる。
  \item 三つ目。
\end{enumerate}

締めの段落。全体で二ページ弱に収まる分量とし、農場の実行時間を
短く保つ。

\end{document}
